\documentclass[12pt]{exam}
\printanswers
%\noprintanswers
\usepackage[T1]{fontenc}        
\usepackage[utf8]{inputenc}   
\usepackage[german]{babel}     
\usepackage[top=2cm,left=2cm,right=2cm,bottom=2cm]{geometry}  
\usepackage{amsmath,amssymb}   
\RequirePackage{amssymb, amsfonts, amsmath, latexsym, verbatim, xspace, setspace}
\usepackage{multicol}         
\usepackage{multirow}     
\usepackage{array}       
\usepackage{ragged2e}    
\usepackage{graphicx}     
\usepackage{pdflscape}    
\usepackage{epstopdf}       
\usepackage{booktabs}       
\usepackage{pdfpages}       
\usepackage[colorlinks=true,urlcolor=magenta,citecolor=red,linkcolor=blue,bookmarks=true]{hyperref}%Links e citações
\usepackage{enumerate}      
\usepackage{enumitem}
\usepackage{color}
\usepackage{tikz,pgfplots} % tikz
\colorgrids
\definecolor{GridColor}{rgb}{0.45, 0.45, 0.45}


\qformat{\textbf{Aufgabe \thequestion\hfill\thepoints}} %%%%%%%% ADDDEDDD %%%%%%


\newlist{partes}{enumerate}{3}
\setlist[partes]{label=(\alph*)}
\newcommand{\parte}{\item}
%---
\newlist{subpartes}{enumerate}{3}
\setlist[subpartes]{label=\roman*)}
\newcommand{\subparte}{\item}
\setlist[enumerate,1]{%(
	leftmargin=*, itemsep=12pt, label={\textbf{\arabic*.)}}}
\newcommand{\dis}{\displaystyle}
\newcommand\answerbox{%%
	\fbox{\rule{2in}{0pt}\rule[-0.1ex]{0pt}{4ex}}}
\pointpoints{Punkt}{Punkte}
\hqword{\textcolor{blue}{Frage:}}
\hpword{Mögliche Punkte:}
\hsword{Erreicht:}
\htword{\textcolor{blue}{Total}}
\renewcommand*\half{.5}
\newcommand{\antwort}[1]{
    {\setlength{\textwidth}{\textwidth}
        \fillwithgrid{#1cm}}
}
\renewcommand{\solutiontitle}{\noindent\textbf{Lösung:}\par\noindent}

% Graphed boxes
\newcommand{\karos}[2]{
  \begin{tikzpicture}
    \draw[step=0.4cm,color=cyan,dotted] (0,0) grid (#1 cm ,#2 cm);
  %  \draw((#1)/2 cm,0)--((#1)/2 cm,/#2)/2cm)
  \end{tikzpicture}
}
%=========================================================

%=========================================================
\newcommand{\thema}{Quadratische Funktionen}
\newcommand{\lehrer}{Philipp Moor}
\newcommand{\klasse}{4f}
\newcommand{\data}{\today}
\newcommand{\dauer}{45 Minuten}
\newcommand{\name}{\bf Name:}
\newcommand{\nota}{Note:}
\newcommand{\datum}{02.12.22}
%=========================================================

%=========================================================
\pagestyle{headandfoot}
\firstpageheader{}{}{}
\runningheader{\datum}{}{Name: \hspace{5cm}}
\runningheadrule
\firstpagefooter{}{}{}
\runningfooter{}{Viel Erfolg!}{Seite \thepage\ von \numpages}
\runningfootrule
%========================================================
\begin{document}
\fontsize{14}{14}\selectfont
%========================================================

%========================================================
\begin{tabular*}{\textwidth}{l @{\extracolsep{\fill}}l @{\extracolsep{6pt}}l}
\textbf{\thema} & \textbf{Note:} & {\answerbox}\\
\textbf{Klasse \klasse} & & \\
\textbf{\lehrer}  & {\bf Name:} &\textbf{\hrulefill}\\[8pt]
%-------------------
\end{tabular*}
%--------------------
\begin{center}
\rule[1ex]{\textwidth}{1pt}
Dauer: \dauer \hspace{9cm}Datum: \datum\\
\vspace{5pt}
\rule[2ex]{\textwidth}{1pt}
\end{center}
%=========================================================

\begin{center}
\textbf{Bewertungstabelle (Wird vom Lehrer ausgefüllt)}\\
\vspace{5pt}
\addpoints
\gradetable[h][questions]\\
\vspace{5pt}
Taschenrechner \textbf{ist erlaubt}, aber jegliche andere Hilfsmittel sind \textbf{nicht erlaubt}.\\
Schreiben Sie alle Blätter mit \textbf{Vor}-und \textbf{Nachnamen} an.
\end{center}
	
\noindent
\rule[1ex]{\textwidth}{1pt}


\newpage

\begin{questions}
%=========================================================
%-------------------TESTFRAGEN
%=========================================================
\noaddpoints
\question[\totalpoints]
\addpoints
\begin{parts}
\part[6] Bestimmen Sie die Funktionsgleichungen $f(x),g(x),h(x)$ anhand folgender Graphen.
(Schreiben Sie die Funktionsgleichungen 
in der Scheitelpunktsform: $f(x)=a(x-u)^2+v$)
\begin{solution}
\begin{align*}
    f(x) &= x^2-4 \\
    g(x) &= \frac{1}{4}(x-2)^2 \\
    h(x) &= -2(x+1)^2+5
\end{align*}
\end{solution}
% evtl weglassen
\part[1] Notieren Sie $k(x)=\frac{1}{8}(x-4)^2+2$ in der Normalform \\ $k(x)=ax^2+bx+c$.
\begin{solution}
\begin{align*}
    k(x) &= \frac{1}{8}(x-4)^2+2 \\
    k(x) &= \frac{1}{8}(x^2-8x+16)+2 \\
    k(x) &= \frac{1}{8}x^2-x+2+2 \\
    k(x) &= \frac{1}{8}x^2-x+4 \\
    %g(x) &= \frac{1}{4}x^2-x+1 \\
    %h(x) &= -2x^2-4x+3
\end{align*}
\end{solution}
\end{parts}

\begin{center}
\begin{tikzpicture}[baseline,scale=0.7]
\begin{axis}[
axis y line=center,
axis x line=middle,
axis equal,
grid=major,
xmax=9,xmin=-7,
ymin=-6,ymax=6,
xlabel=$x$,ylabel=$y$,
xtick={-10,-9,...,10},
ytick={-10,-9,...,10},
width=15cm,
anchor=center,
]

\addplot[samples=250,mark=none,color=red,domain=-4:4]{x^2-4} node [pos=0.7, above right] {$f(x)$};
\addplot[samples=250,mark=none,color=blue,domain=-6:9]{0.25*(x-2)^2} node [pos=0.7, below right] {$g(x)$};
\addplot[samples=250,mark=none,color=brown,domain=-4:2]{-2*(x+1)^2+5} node [pos=0.5, above] {$h(x)$};
\end{axis}
\end{tikzpicture}
\end{center}

\ifprintanswers\else
    \vspace{5pt}
    \antwort{11}
\fi
\newpage

\noaddpoints
\question[\totalpoints]
\addpoints
Gegeben seien die Funktionen \[f(x)=x^2-6x+8 \text{ , } \quad g(x) = x^2-7x+6\]
\begin{parts}
\setlength{\columnsep}{30pt}
\part[2] Skizzieren sie den Graphen der Funktionen $f(x)$ in folgendem \\ Koordinatensystem:
(Kennzeichnen Sie die berechneten Punkte)
\begin{solution}
\begin{center}
\begin{tikzpicture}[baseline,scale=0.5]
\begin{axis}[
axis y line=center,
axis x line=middle,
axis equal,
grid=major,
xmax=6,xmin=-1,
ymin=-3,ymax=5,
xlabel=$x$,ylabel=$y$,
xlabel style={at={(ticklabel* cs:1)},anchor=north west},
ylabel style={at={(ticklabel* cs:1)},anchor=south west},
xtick={-10,-9,...,10},
ytick={-10,-9,...,10},
width=15cm,
anchor=center,
]
\addplot[samples=250,mark=none,color=red,domain=0:6]{x^2-6*x+8} node [pos=0.7, above right] {$f(x)$};
\end{axis}
\end{tikzpicture}
\end{center}
\end{solution}
\ifprintanswers\else
    \begin{center}
    \begin{tikzpicture}[baseline,scale=0.5]
    \begin{axis}[
    axis y line=center,
    axis x line=middle,
    axis equal,
    grid=major,
    xmax=6,xmin=-1,
    ymin=-3,ymax=5,
    xlabel=$x$,ylabel=$y$,
    xlabel style={at={(ticklabel* cs:1)},anchor=north west},
    ylabel style={at={(ticklabel* cs:1)},anchor=south west},
    xtick={-10,-9,...,10},
    ytick={-10,-9,...,10},
    width=15cm,
    anchor=center,
    ]
    %\addplot[samples=250,mark=none,color=red,domain=0:6]{x^2-6*x+8} node [pos=0.7, above right] {$f(x)$};
    \end{axis}
    \end{tikzpicture}
    \end{center}
\fi
\part[2] Berechnen Sie die Nullstellen der Funktion $g(x)$ ohne TR.
\begin{solution}
    \begin{align*}
        g(x)&=x^2-7x+6 \\
        g(x)&=(x-6)(x-1) \\
        x_1 &= 6 \quad x_2 = 1
    \end{align*}
\end{solution}
\part[2] Berechnen sie den Scheitelpunkt der Funktion $g(x)$ anhand der Nullstellen. (ohne TR)
\begin{solution}
Nullstellen bei $x_1=6$ und $x_2=1$ $\implies$ $x$-Wert des Scheitelpunkts bei $3.5$
\begin{align*}
    g(3.5)&=(3.5)^2-7\cdot3.5+6 \\
    &= -6.25
\end{align*}
$\implies$ $y$-Wert bei -6.25 $\implies$ Scheitelpunkt $=(3.5,-6.25)$
\end{solution}
\part[2] Bringen Sie die Funktion $g(x)$ in Scheitelpunktsform. \\ 
($g(x)=a(x-u)^2+v$)
\begin{solution}
Scheitelpunkt bei $(3.5,-6.25) \implies$
\begin{align*}
    g(x)&=a(x-3.5)^2-6.25 \\
    0&=a(1-3.5)^2-6.25 \\
    a(1-3.5)^2&=6.25 \\
    6.25a&=6.25 \\
    a&=\frac{6.25}{6.25} \\
    &= 1
\end{align*}
$\implies$ $g(x)=(x-3.5)^2-6.25$
\end{solution}
\end{parts}
\ifprintanswers\else
    \vspace{5pt}
    \antwort{10}
\fi
\newpage

\noaddpoints
\question[\totalpoints]
\addpoints
Veränderungen von Funktionsgleichungen:\\
(Notieren sie bei a),b) jeweils die veränderte Funktionsgleichung 
als $f'(x),g'(x)$)
\begin{parts}
\setlength{\columnsep}{30pt}
\part[2] Verschieben sie $f(x)=3x^2-4$ \\ um 3 Einheiten nach Rechts und um 5 Einheiten nach Oben.
\begin{solution}
\begin{align*}
    f'(x)=3(x-3)^2+1
\end{align*}
\end{solution}
\part[2] Verschieben sie $g(x)=-(x-3)^2+3$ \\ um 4 Einheiten nach Rechts und um 2 Einheiten nach Unten.
\begin{solution}
\begin{align*}
    g'(x)=-(x+1)^2-1
\end{align*}
\end{solution}
%\part Verschieben sie $h(x)=-\frac{1}{2}x^2-4x+3$ um 3 Einheiten nach Links und um 5 Einheiten nach Unten.
\part[2] Um wieviele Einheiten nach Rechts/Links resp. nach Oben/Unten wurde verschoben wenn $h(x)$ zu $h'(x)$ wird ?
\begin{center}
    $h(x)=(x-\frac{1}{2})^2-3 \quad \rightarrow \quad h'(x)=(x+4)^2+9$
\end{center}
\begin{solution}
Um 4.5 Einheiten nach Links und um 12 nach Oben.
\end{solution}
\end{parts}

\ifprintanswers\else
    \newpage
    \antwort{24}
\fi
\newpage
\noaddpoints
\question[\totalpoints]
\addpoints
% skizze und gebe gleichung an...

Ein Bauer baut auf einer Flachen Wiese ein rechteckiges Gehege für seine Hühner. Er hat Zaun in der länge von $20m$ gekauft und will den ganzen Zaun verbauen (der Umfang des entstehenden Rechtecks soll $20m$ messen). Es gilt:
\begin{align}
    2a+2b&=20m  \\
    \text{Fläche } A&=a\cdot b
\end{align}
wenn $a$ und $b$ die Seiten des Rechtecks sind.
\begin{parts}
\part[1] Formulieren Sie die Gleichung (1) nach $a$ um ($a=\ldots$).
\part[1] Setzen Sie das berechnete $a$ in die Gleichung (2) ein ($A=\ldots\quad\cdot b$).
\part[2] Skizzieren sie die Funktion aus (b) ins folgende Koordinatensystem: \\ (falls Sie (a) und (b) nicht berechnen konnten: Skizzieren Sie $A=-\frac{1}{2}b^2+5b$)
\ifprintanswers\else
    \begin{center}
    \begin{tikzpicture}[baseline,scale=0.8]
    \begin{axis}[
    axis y line=center,
    axis x line=middle,
    %axis equal,
    grid=major,
    xmin=0,
    xmax=11,
    ymin=0,
    ymax=26,
    xlabel=$b$,ylabel=$A$,
    xlabel style={at={(ticklabel* cs:1)},anchor=north west},
    ylabel style={at={(ticklabel* cs:1)},anchor=south west},
    xtick={-1,0,...,11},
    ytick={0,1,...,25},
    width=15cm,
    anchor=center,
    ]
    %\addplot[samples=250,mark=none,color=red,domain=0:10]{-x^2+10*x} node [pos=0.7, above right] {$f(x)$};
    %\addplot[samples=250,mark=none,color=red,domain=0:10]{-(1/2)*x^2+5*x} node [pos=0.7, above right] {$f(x)$};
    \end{axis}
    \end{tikzpicture}
    \end{center}
\fi
\part[3] Wie sind die Seitenlängen $a$ und $b$ des Rechtecks zu wählen wenn das Gehege $16m^2$ Fläche haben soll?
\part(1 Bonuspunkt) Bei welchen Seitenlängen $a$ und $b$ is das Gehege maximal gross?
\end{parts}
\ifprintanswers\else
    \newpage
    \antwort{24}
\fi
\newpage
\noaddpoints
\question[\totalpoints]
\addpoints
Gegeben seien die folgenden Funktionen:
\begin{align*}
    f(x)= 1.75x^2-3x-2 \text{ , } \quad g(x) = -x^2+2.5x-2
\end{align*}
\begin{parts}
\part[1] Berechnen sie die Nullstellen von $f(x)$ und $g(x)$ (TR möglich).
\part[1] Berechnen sie die Scheitelpunkte von $f(x)$ und $g(x)$ (TR möglich).
%\part Skizzieren Sie die beiden Graphen $f(x)$ und $g(x)$ ins gleiche Koordinatensystem (\textbf{Achten sie auf korrekte Beschriftung der Achsen})
\part[4] Berechnen Sie die $x$- und $y$-Koordinaten der Schnittpunkte von $f(x)$ mit $g(x)$.
\end{parts}
\ifprintanswers\else
    \vspace{5pt}
    \antwort{18}
\fi
\newpage


\noaddpoints
\question[\totalpoints]
\addpoints
Gegeben seien die folgenden Funktionen:
\begin{align*}
    f(x)= 2x^2+8x+7 \text{ , } \quad g(x) = 5x^2+4x+3
\end{align*}
\begin{parts}
\part[\half] Berechnen Sie mittels TR die Nullstellen von $f(x)$.
\part[2] Berechnen Sie von Hand mittels Lösungsformel (Mitternachtsformel) die Nullstellen von $f(x)$.
\part[1] Argumentieren Sie in ein bis zwei Sätzen, welcher Vorteil die Berechnung der Nullstellen ohne TR bieten kann.
\part[1] Weshalb können die Nullstellen von $g(x)$ nicht \\
berechnet werden ? (Argumentieren Sie in ein bis zwei Sätzen)
\end{parts}

\ifprintanswers\else
    \vspace{5pt}
    \antwort{17}
\fi

\addpoints
\end{questions}
\end{document}